\subsection{Objetivo del refinamiento}

Checkpoint~04D.1 no vuelve a realizar una búsqueda genérica de modelos. Toma como evidencia
previa que LocalRidge y Graph sobrevivieron 04B, reutiliza exactamente las 34 máscaras
congeladas y explora únicamente grados de libertad relevantes para estas dos familias.

El grid principal contiene 626 métodos predeclarados:
\begin{itemize}
\item 300 variantes de grafo;
\item 324 variantes de Local Ridge;
\item dos anchors fijos.
\end{itemize}
La gran cantidad de variantes no se interpreta como 626 hipótesis independientes; comparten
estructura y muchos hiperparámetros adyacentes. Por eso se estudian superficies de estabilidad
y luego se usa selección leave-one-split-out.

\subsection{Topologías de distancia}

Las topologías de distancia combinan geografía y embedding agronómico:
\[
d^{(\eta)}_{ij}
=
\eta\,\widetilde d^{geo}_{ij}
+
(1-\eta)\,\widetilde d^{agro}_{ij},
\]
con variantes denominadas Geo, GeoAgro90, GeoAgro75, GeoAgro50 y GeoAgro25. En la
nomenclatura, GeoAgro25 corresponde a 25\% geografía y 75\% componente agronómica
normalizada.

Local Ridge se prueba sobre C1 agronomic y C4 all deterministic, con topologías Geo,
GeoAgro75 y GeoAgro50. La búsqueda usa
\[
k\in\{12,16,20,24,30,40\},
\]
\[
\alpha\in\{30,100,300\},
\qquad
p\in\{0.5,1,2\},
\]
donde $p$ es la potencia de distancia en
\[
w_{qi}=\frac{1}{(d_{qi}+10^{-6})^p}.
\]

\subsection{Local Ridge query-specific}

Para cada objetivo $q$ se seleccionan únicamente etiquetas visibles:
\[
\mathcal N_k(q)
=
\operatorname*{arg\,sort}_{i\in\Train^\star}
d_{qi}[1:k].
\]
Entonces
\begin{equation}
\hat{\bm\beta}_q
=
\arg\min_{\bm\beta}
\sum_{i\in\mathcal N_k(q)}
w_{qi}
(y_i-x_i^\top\bm\beta)^2
+
\alpha\|\bm\beta\|_2^2.
\label{eq:cp04d-localridge}
\end{equation}

La regularización es esencial: incluso con $k=24$, C4 tiene una dimensión mucho mayor que el
tamaño del vecindario. Ridge no intenta identificar un vector causal de coeficientes; actúa
como interpolador lineal regularizado en el subespacio local. La estabilidad proviene de
$\alpha$ elevado respecto a un OLS local singular.

\subsection{Búsqueda de grafos}

Los grafos directos varían:
\[
k\in\{4,6,8,10,12\},\qquad
\lambda\in\{0.25,0.5,1,2,4,8\}
\]
sobre cinco topologías. La solución sigue \cref{eq:graph-laplacian}. Valores mayores de
$\lambda$ penalizan más la variación de la función sobre aristas:
\[
\bm f^\top L\bm f
=
\frac{1}{2}\sum_{i,j}w_{ij}(f_i-f_j)^2.
\]
Por tanto, el modelo asume smoothness local del rendimiento en el grafo $X$.

\subsection{Corrección de residuos globales}

Además del grafo directo se probó un modelo de corrección de residuos. Primero se ajusta PLS4
sobre etiquetas visibles y se generan residuos cross-fitted dentro de ese mismo conjunto:
\[
r_i^{CF}
=
y_i-\hat y_i^{PLS,-fold(i)}.
\]
Se ajusta después un campo suave sobre el grafo:
\[
\hat r_q
=
\operatorname{GraphPropagate}(r^{CF}),
\]
y se predice
\[
\hat y_q
=
\hat y_q^{PLS,full}+\hat r_q.
\]
El uso de residuos cross-fitted es crítico. Si se propagaran residuos in-sample de un PLS
ajustado sobre los mismos puntos, su magnitud y estructura estarían contaminadas por fit
optimista.

\subsection{Ranking target-matched}

\begin{table}[H]
\centering
\caption{Mejores configuraciones del refinamiento 04D.1.}
\label{tab:cp04d-ranking}
\small
\begin{tabular}{lrrr}
\toprule
Método & RMSE medio & mediana & peor split \\
\midrule
C4 Geo $k=24,\alpha=30,p=1$ & \textbf{0.487845} & 0.468742 & 0.694700 \\
C1 Geo $k=30,\alpha=30,p=1$ & 0.488302 & 0.4694 & 0.6921 \\
C1 Geo $k=24,\alpha=30,p=1$ & 0.4885 & 0.4684 & 0.6953 \\
C4 Geo $k=30,\alpha=30,p=1$ & 0.4887 & 0.4715 & 0.6925 \\
C4 GeoAgro75 $k=24,\alpha=30,p=1$ & 0.4892 & 0.4718 & 0.6940 \\
\bottomrule
\end{tabular}
\end{table}

La proximidad de varias configuraciones alrededor de $k=20$--30 y $\alpha=30$ es más
importante que el orden exacto de las primeras filas: existe una cuenca estable, no un spike
aislado. Esta observación aumenta la credibilidad del finalista.

\safefigure{../../reports/checkpoint_04d1/figures/local_k_profile.png}
{Perfil de desempeño de Local Ridge al variar el tamaño del vecindario.}
{fig:cp04d-local-k}

\subsection{El mejor modelo robusto no coincide con el mejor target-matched}

El finalista de grafo más robusto es
\artifact{GraphDirect_GeoAgro25_k6_lam8}:
\[
\begin{array}{lr}
\text{target-matched} & 0.494881\\
\text{state-random} & 0.467936\\
\text{state-stratified} & 0.479161\\
\text{municipality-grouped} & 0.527883.
\end{array}
\]

El mejor Local Ridge obtiene 0.487845 target-matched, pero municipality-grouped cerca de
0.687. La diferencia expresa dos objetivos estadísticos:
\begin{itemize}
\item Local Ridge aprovecha intensamente etiquetas geográficamente cercanas, condición que
el matching indica que se parece al test real;
\item GraphDirect conserva mejor comportamiento cuando se elimina un municipio completo.
\end{itemize}

No existe contradicción en conservar ambos como finalistas: uno maximiza la simulación primaria
del test fijo; el otro es un hedge de extrapolación espacial.

\safefigure{../../reports/checkpoint_04d1/figures/refinement_ranking.png}
{Ranking del search local/grafo sobre pseudo-competiciones target-matched.}
{fig:cp04d-refinement}

\subsection{Graph residual versus graph directo}

Las variantes de corrección de residuos PLS quedan alrededor de 0.537 target-matched y
0.679--0.684 grouped, claramente por debajo del grafo directo. Esto sugiere que, en este
problema, la estructura espacial no actúa sólo como una pequeña corrección de un anchor global:
modelar directamente la función sobre el grafo es más competitivo.

\subsection{Leave-one-split-out para selección de método}

Para cada split $h$:
\[
\hat m_{-h}
=
\arg\min_m
\frac{1}{15}
\sum_{s\neq h}\RMSE_{s,m},
\]
y luego se calcula $\RMSE_{h,\hat m_{-h}}$. El resultado es:
\[
\overline{\RMSE}_{LOSO,\mathrm{method}}=0.488366,
\qquad
\RMSE_{\mathrm{pooled}}=0.496137.
\]
Quince de dieciséis selecciones eligieron el C4 Geo $k=24,\alpha=30,p=1$ y una eligió C1
Geo $k=30,\alpha=30,p=1$. Esta concentración confirma que el optimum no depende de un solo
split.

\subsection{Routing por soporte y caveat posterior}

04D.1 también ejecutó selección LOSO por tier y reportó:
\[
\overline{\RMSE}_{LOSO,\mathrm{tier}}=0.485815,
\qquad
\RMSE_{\mathrm{pooled}}=0.494213.
\]
Numéricamente es mejor que la regla única. Sin embargo, una auditoría posterior detectó una
asimetría importante: el universo del selector LOSO de tiers podía incluir los 626 métodos
predeclarados, mientras la propuesta actual para las 59 parcelas se construía sobre un
subconjunto de finalistas. Una ganancia de aproximadamente 0.0026 no justifica ignorar esa
diferencia de candidate universe.

Además, los 16 pseudo-splits reutilizan parcelas; no son 16 datasets independientes. La
selección split-excluded remueve optimismo de mismo split, pero no convierte los resultados
en replicaciones estadísticamente independientes.

\begin{decision}
El routing por soporte se conserva como evidencia exploratoria, no como regla final. La
solución final debe usar un candidate universe explícito y estable o elegir una regla única
suficientemente robusta.
\end{decision}

\subsection{Candidato para las 59 parcelas}

04D.1 materializa predicciones candidatas, incluyendo Local C4 y GraphDirect. Estos valores no
se promueven todavía a entrega: sirven para medir discrepancia entre inductive biases y para
que 04E/04C evalúen si evidencia externa o global puede reducir error antes del congelamiento
de 04F.
